% Problem. The paper's own formulation: families with a vocabulary, rule lists and their meaning on an
% instance, certification, margin, and the problem statement. Builds on sec:background.

\section{Problem Definition}
\label{sec:problem}

In this section, we define the families of instances our framework controls (\autoref{sec:families}), the rule lists it writes and their meaning on an instance (\autoref{sec:rules}), and when a rule list is certified, which leads to the problem we solve (\autoref{sec:certified}).

\subsection{Families and Their Vocabulary}
\label{sec:families}

A \textit{family} $\mathcal{F}$ is a set of instances, each given by an MDP $M_I$ (\autoref{sec:bg-mdp}).
The instances share a set of \textit{policy actions} $A_\pi$, which the controller chooses, and a set $\Phi$ of requirements, whose thresholds $\theta^I$ may differ between instances.
Any other choice in the model, e.g., a move of the environment, is left open and resolved by the worst case.
The instances also share a set $X$ of \textit{observable} variables with the same names and types, and the controller only sees the values $\mathit{obs}(s)$ of these variables in a state $s$.

What sets one instance apart, e.g., the positions of its goals, is not part of the state.
To let one controller speak about every instance, the domain declares a \textit{vocabulary} with three kinds of terms.
\textit{Constants} are named values that each instance fixes, e.g., the row and column of each goal, or the length of a Pac-Man game.
\textit{State features} are expressions over the observable variables and the constants, e.g., the number of moves left in a Pac-Man game.
\textit{Action features} have one expression per policy action, e.g., \texttt{closer} holds for a move if it brings the agent closer to the next goal along a shortest path.
In the running example, the observable variables are the agent's cell, the goal flags and the position of the moving obstacle, and \autoref{fig:intro-rules} uses the action features \texttt{wall}, \texttt{closer}, \texttt{risky}, \texttt{obstacle}, \texttt{early} and \texttt{patrol}.
\subsection{Rule Lists}
\label{sec:rules}

A \textit{rule list} $R = \langle (c_1, a_1), \ldots, (c_n, a_n) \rangle$ is an ordered list of rules, each with a condition $c_j$ and an action $a_j \in A_\pi \cup \{\texttt{any}\}$.
Conditions are Boolean combinations of comparisons over the observable variables and the vocabulary.
To keep a rule list meaningful on every instance, the only numbers a condition may contain are 0 and 1, so conditions name instance values through constants instead of writing them out\todo{Pac-Man allows other numbers on variables whose range is the same on every instance (rules.shared\_literals); say so here or in the evaluation}.

On an instance $I$, a rule $(c_j, a)$ with $a \in A_\pi$ \textit{allows} $a$ in state $s$ if $c_j$ holds in $\mathit{obs}(s)$, where every action feature in $c_j$ is evaluated for $a$.
A rule $(c_j, \texttt{any})$ allows each action $b \in A_\pi$ for which $c_j$ holds with its action features evaluated for $b$.
The first rule that allows some action in $s$ \textit{decides} $s$, so earlier rules take precedence and overlapping conditions never conflict.
The rule list allows the actions its deciding rule allows, and every policy action in a state that no rule decides, which we call \textit{uncovered}.
For example, rule~2 of \autoref{fig:intro-rules} allows every move that does not hit a wall, brings the agent closer to the next goal, cannot slip onto a hazard, and neither meets the moving obstacle nor enters a goal out of turn.
Where no move satisfies all of this, a later, weaker rule decides, and the final rule \texttt{true} leaves no state uncovered.

Restricting $M_I$ to the actions $R$ allows yields the MDP $M_I^R$, i.e., a permissive strategy in the sense of \autoref{sec:bg-pmc}.
We build $M_I^R$ by adding one PRISM module that synchronizes on the policy actions and blocks every action the rules do not allow.
Its guard for action $a$ is
\begin{equation}
\label{eq:guard}
\bigvee_{j=1}^{n} \Big( c_j[a] \wedge \bigwedge_{k < j} \neg m_k \Big) \;\vee\; \bigwedge_{k=1}^{n} \neg m_k,
\end{equation}
where the disjunction ranges over the rules whose action is $a$ or \texttt{any}, $c_j[a]$ is $c_j$ with its action features evaluated for $a$, and $m_k$ holds where rule $k$ allows some action.
In words, $a$ is enabled where the deciding rule allows it, or where no rule decides.

\subsection{Certified Rule Lists}
\label{sec:certified}

The worst case $w_\varphi^I(R)$ and the best case $b_\varphi^I(R)$ of a requirement $\varphi \in \Phi$ on instance $I$ are its worst and best case on $M_I^R$, taken over all its strategies.
A rule list $R$ is \textit{certified} on $I$ when $w_\varphi^I(R)$ meets the threshold $\theta^I$ of every requirement $\varphi$.
By \autoref{sec:bg-pmc}, every controller that only takes actions $R$ allows then meets every requirement on $I$, however it breaks ties among the actions of an \texttt{any} rule, whatever it does in uncovered states, and even if it uses memory or randomization.
The worst case is also conservative in what it observes, since the strategies of $M_I^R$ see the full state while the rules see only the observable variables.
If a rule list allows a single action in every reachable state, its worst and best cases coincide and equal the value of the one controller it defines.

The \textit{margin} of $R$ on $I$ for $\varphi$ is how far its worst case lies beyond the threshold, i.e., $w_\varphi^I(R) - \theta^I$ for a lower bound and $\theta^I - w_\varphi^I(R)$ for an upper bound.
$R$ is certified on $I$ exactly when its smallest margin over all requirements is non-negative, and the smallest margin also ranks rule lists that are not yet certified by how close they come.

\textbf{Problem Definition:}
Given a family $\mathcal{F}$ with requirements $\Phi$ and a vocabulary, and a sequence of instances $I_1, I_2, \ldots \in \mathcal{F}$ that the system meets at run time, the problem is to provide, for each $I_t$, a rule list certified on $I_t$ before it controls $I_t$.
Rule lists should be reused across instances wherever they are certified, since checking a rule list on an instance is cheaper than writing a new one.
